\chapter{Project Practice and Task Tracking}
\label{ch:project-practice}

Chapter~\ref{ch:design-philosophy} covered the code-level house style CNA
enforces on every file. This chapter covers the surrounding \emph{project} practice --- how
work is planned, tracked, and handed off between sessions --- because reading CNA's own
repository well enough to contribute to it, or even to just trust a specific claim in it,
means knowing this layer exists and how to read it.

\section{Three tiers of tracking document}

CNA's root directory carries three distinct kinds of tracking file, each with a different job,
and confusing one for another is a real way to misread the project's own state.

\textbf{\texttt{NEXT.md}} is a \emph{session handoff} document --- written for whichever
agent or engineer picks up work next, not a historical record. It is explicitly scoped to
whichever development branch is currently active, and is candid about its own staleness in a
way worth quoting directly, since it is a good model of the honesty this book has tried to
carry throughout: one recent revision opens by telling its own reader that everything below a
certain divider is ``stale, D3D9-branch-scoped content\ldots that ended up merged into this
file's history\ldots not about this branch and should not be read as current status here.''
Rather than deleting the stale material, it is left in place as a historical record, with an
explicit warning attached. \texttt{NEXT.md} also carries genuinely operational notes with no
equivalent anywhere else in the project --- for instance, running graphics/window tests
against a dedicated Xvfb display rather than a real one, capping build parallelism well below
the machine's full core count and pausing new heavy work above a specific CPU temperature
threshold when the development machine is shared concurrently with other agents, and a
reminder that \texttt{ctest} changes each test's working directory, which can silently break a
test relying on a path relative to the repository root.

\textbf{\texttt{AUDIT.md}} is a \emph{systematic, per-class comparison} between FNA (as
reference) and CNA (as implementation), using a three-value legend --- done, in progress, not
yet audited --- across every class and enum in every XNA namespace. Unlike \texttt{NEXT.md},
this is meant to be read as a durable, cumulative record of what has actually been checked,
not a snapshot of current work.

\textbf{\texttt{CHECKLIST.md}} is the per-file porting checklist already introduced in
Chapter~\ref{ch:design-philosophy} --- SPDX headers, the \texttt{NOXNA} include requirement,
line-by-line verification against FNA, mandatory \texttt{GetTypeName()} overrides, and full
test coverage including out-ref overloads tested separately from their value-returning
counterparts. It also maintains the project's own table of \emph{accepted} C\texttt{++}
deviations from FNA --- for instance, the inconsistency already named in
Chapter~\ref{ch:math-core-types} between \cnaclass{BoundingBox} throwing
\texttt{System::ArgumentOutOfRangeException} for an out-of-range index versus
\cnaclass{Input::Touch::TouchCollection} throwing \texttt{std::out\_of\_range} for the
identical situation --- flagged explicitly as a known, project-wide inconsistency rather than
silently tolerated or silently ``fixed'' to look more consistent than the codebase's actual
history.

\section{One plan file per subsystem}

Beyond these three root-level files, CNA maintains a separate \texttt{plan\_<subsystem>.md}
file for each major feature area under active or completed development --- at the time of
writing, twenty of them, covering areas as different in scale as \texttt{plan\_ascii.md} (the
ASCII backend, Chapter~\ref{ch:canvas-ascii-backends}) and \texttt{plan\_cnj.md}
(the glTF import pipeline). Each is organized into numbered phases and tasks --- the same
\texttt{Task~123} citation style this book has used throughout both volumes whenever
attributing a specific bug or fix to its source. This is the mechanism that makes a claim
like ``Task~868, fixed 2026-07-09'' (Chapter~\ref{ch:vulkan-backend}) a genuinely checkable
one rather than an unsourced assertion: the task number names an exact entry in an exact
plan file, not a vague reference to ``recent work.''

\section{Status symbols, read consistently}

Across \texttt{NEXT.md} and the \texttt{plan\_*.md} files, a small set of symbols recur with
consistent meaning: a hollow marker for work not yet started, a half-filled marker for a
partially complete item, and a checkmark for done --- the same three-state shape
\texttt{AUDIT.md} uses with different symbols. When this book's own chapters describe a
feature as ``open'' or ``BLOCKED'' rather than simply absent, that distinction is drawn
directly from this same convention: an open item is ordinary unstarted or in-progress work; a
BLOCKED item, as Chapter~\ref{ch:sdl-renderer} covered for \texttt{SDL\_RENDERER}'s two
BLOCKED rows, specifically means the work itself is understood but genuinely cannot proceed
without a decision only the project owner can make.

\section{Doxygen \texttt{@note Status:} tags: a secondary, human-readable hint}

\texttt{sharp-runtime} (Chapter~\ref{ch:sharp-runtime-overview}) layers one further,
source-level convention on top of all of this: a \texttt{@note Status: Todo|Stub|Partial|%
Implemented|Verified} Doxygen tag directly on individual classes and methods. Its own
documentation is explicit that this tag is a \emph{secondary}, human-readable hint, not a
source of truth --- the authoritative record of whether a given .NET type counts as ``ported''
lives in \texttt{plan.sqlite3}'s own \texttt{task} table (Chapter~\ref{ch:sharp-runtime-overview}
covers this database in full), and the two are not mechanically kept in sync with each other.
A reader trusting only the source comment risks trusting a claim the project's own
authoritative tracking database might already have superseded.

\section{Why this layer is worth knowing, not just tolerating}

None of this is unusual in isolation --- plenty of large projects keep a TODO file and a set
of per-feature planning documents. What is worth taking away from this chapter is the
discipline threading through all of it: every one of these documents is written to be
distrusted a little, and to tell you exactly how to check it against something more
authoritative --- \texttt{NEXT.md} telling you which of its own sections are stale,
\texttt{AUDIT.md} tracking currency per class rather than claiming a single global percentage,
\texttt{CHECKLIST.md} recording known inconsistencies instead of hiding them, and
\texttt{sharp-runtime}'s own Doxygen tags explicitly deferring to a database rather than
claiming to be the final word themselves. This is the same discipline
Chapter~\ref{ch:what-is-cna} named as this project's most distinctive trait in its very first
pages, and it is worth recognizing it here again, now grounded in the specific documents that
actually carry it.
